\documentclass{article}

% -------------------------------------------------------
% FORMAT
% -------------------------------------------------------
% Input encoding and language.
\usepackage[spanish, es-tabla]{babel}
\usepackage[utf8]{inputenc}
% Refer to p.92 of "The not so short introduction to Latex" 
% by Tobias Oetiker to understand what role have this packages.
\usepackage{lmodern}
\usepackage[T1]{fontenc}
\usepackage{textcomp}

% Provide a high-level, flexible interface for creating document-level commands and environments.
\usepackage{xparse}

% Verbatim text formatting, including headers.
\usepackage{fancyvrb}
% Boxes management.
\usepackage{fancybox}
% Headers management.
\usepackage{fancyhdr}
% First you need to clear the default layout.
\fancyhead{}
\fancyfoot{}
\fancyhead[RE]{\rightmark}
\fancyhead[LO]{\rightmark}
\fancyfoot[C]{\thepage}
% Remove decorative lines.
\renewcommand{\headrulewidth}{0pt}
\renewcommand{\footrulewidth}{0pt}
\pagestyle{fancy}

% Package for underlining, adding color to text, and highlighting text in color.
% Text highlighting package. It solves the 'fullbox' errors when using \colorbox{declared-color}{text}.
% It doesn't support letters with accents, so they must be included using the '\'' syntax.
% Usage: \hl{text}.
\usepackage{soul}
% Sets the 'soul' package to higlight text using color red.
\sethlcolor{red}

% -------------------------------------
% Easy quoting.
\usepackage[style=american]{csquotes}


% ------------------------------------------------------------------------
% MATHS
% ------------------------------------------------------------------------
\usepackage{amsmath,amsfonts,amsthm,amssymb,amsbsy}
% Equation numbering format by section. E.g., section x, equation y: (x.y)
\numberwithin{equation}{section}
\numberwithin{figure}{section}
\usepackage{mathtools}
% Useful for arrays of equations.
\usepackage[retainorgcmds]{IEEEtrantools}

% ------------------------------------------
% Setting Up Automatic Scaling Brackets.
% ------------------------------------------
% When you write equations with tall elements like fractions or matrices, standard parentheses ( and ) stay small and look awkward. While typing \left( and \right) fixes this, doing it every single time gets exhausting and clutters your code.
% Define custom paired delimiters
\DeclarePairedDelimiter\paren{(}{)}       % Parentheses
\DeclarePairedDelimiter\sqbrac{[}{]}      % Square brackets
\DeclarePairedDelimiter\curly{\{}{\}}     % Curly braces
\DeclarePairedDelimiter\abs{\lvert}{\rvert} % Absolute value
\DeclarePairedDelimiter\norm{\lVert}{\rVert} % Norm

% ----------------------------
% Imaginary unit.
% The international standard for mathematical typesetting (ISO 80000-2) dictates that mathematical constants—like the base of the natural logarithm ($e$), pi ($\pi$), and the imaginary unit ($i$ or $j$)—should be typeset in an upright (Roman) font to clearly distinguish them from variables. (Notice how the 'i' stands straight up when using it, making it instantly recognizable as a constant rather than a variable exponent).
% ----------------------------
\newcommand{\iu}{\mathrm{i}} % Imaginary unit for math/physics
\newcommand{\ju}{\mathrm{j}} % Imaginary unit for engineering
% ----------------------------
% Generic counter
% ----------------------------
\newcounter{hipotesis_ct}
\setcounter{hipotesis_ct}{1}

% ----------------------------
% Derivatives operators.
\usepackage{derivative}

% ----------------------------
% Set operators.
% ----------------------------
% We create a dummy command \given that we will redefine inside our \set command
\providecommand\given{} 
% The 'X' allows us to inject formatting inside the delimiters
\DeclarePairedDelimiterX\set[1]\{\}{
	% This makes \given print a vertical bar that scales with the outer braces
	\renewcommand\given{\;\delimsize\vert\;}
	#1
}

% ----------------------------
% Commands to write vectors.
% ----------------------------
\usepackage[c]{esvect}
\newcommand{\vecbf}[1]{\mathbf{#1}}


% ----------------------------
% Commands to write vectors.
% ----------------------------
\usepackage[c]{esvect} % Keep this ONLY if you actually use \vv{x} for arrows somewhere
\usepackage{bm}        % The modern bold math package

% Core semantic definitions
\newcommand{\mVec}[1]{\bm{#1}}       % Bold italic vectors (works on Greek too!)
\newcommand{\Mat}[1]{\bm{#1}}        % Matrices (often same as vectors, but good to separate semantically)
\newcommand{\quat}[1]{\mathbf{#1}}   % Quaternions (Standard is upright bold)

% Note: The \dot goes OUTSIDE the \mVec so the dots don't get bolded.
% Vector decorations
\newcommand{\vechat}[1]{\hat{\mVec{#1}}}
\newcommand{\vecunit}[1]{\breve{\mVec{#1}}}
\newcommand{\vecdot}[1]{\dot{\mVec{#1}}}      % Velocity
\newcommand{\vecddot}[1]{\ddot{\mVec{#1}}}    % Acceleration
\newcommand{\vecdddot}[1]{\dddot{\mVec{#1}}}  % Jerk
\newcommand{\vecprod}[2]{#1 \times #2}

% Quaternion operations
\newcommand{\quatComp}[2]{#1 \circ #2}	% Composition.
\newcommand{\quatConj}[1]{#1^{\ast}}	% Conjugate.
\newcommand{\quatRot}[2]{\Re(#1, #2)} 	% Rotation.
% ----------------------------
% Estimators.
% Dynamically stretches the hat to perfectly cover whatever is inside of it.
\newcommand{\estimated}[1]{\widehat{#1}}

% ----------------------------
% Command for equations definitions.
% Requires the amssymb package
\newcommand{\defineEq}{\triangleq}

% ----------------------------
% Système international d'unités
\usepackage{siunitx}

% ---------------------------------
% Automatic matrix transposition.
% ---------------------------------
% This snippet uses LaTeX3 (expl3) programming to create a highly useful new math environment called bmatrixT. In short, it automatically transposes a matrix. You type the matrix normally, and when the document compiles, the code swaps the rows and columns and wraps the result in standard square brackets (a bmatrix).
% It also works with pmatrix (parentheses) and vmatrix (vertical bars).
%
% How it works (using the bmatrixT environment as example):
%
% Capturing the Content: It uses the environ package to grab everything you type inside \begin{bmatrixT} ... \end{bmatrixT} and saves it as a single chunk of text (\BODY).
%
% Deconstructing the Matrix: It splits the input at every line break (\\) to figure out the rows.
%
% It then splits each row at every ampersand (&) to separate the individual column items.
%
% It stores every single item in a temporary property list using its coordinates (Row, Column) as a tag.
%
% Rebuilding the Matrix (The Transpose): It uses nested loops to rebuild the matrix, but it reverses the order. It loops through the columns first, and the rows second.
%
% This effectively swaps the X and Y coordinates of every item.
%
% Formatting the Output: Finally, it takes the newly swapped items, adds the & and \\ symbols back in the correct spots, and feeds the whole thing into a standard bmatrix environment.
\usepackage{environ}

\ExplSyntaxOn
% Declare variables
\int_new:N \l_marine_transpose_row_int
\int_new:N \l_marine_transpose_col_int
\seq_new:N \l_marine_transpose_rows_seq
\seq_new:N \l_marine_transpose_arow_seq
\prop_new:N \l_marine_transpose_matrix_prop
\tl_new:N \l_marine_transpose_last_tl
\tl_new:N \l_marine_transpose_body_tl

% 1. The function takes TWO arguments (:nn)
% #1 is the environment type (bmatrix, pmatrix, etc.)
% #2 is the matrix content
\cs_new_protected:Nn \marine_transpose:nn
{
	\seq_set_split:Nnn \l_marine_transpose_rows_seq { \\ } { #2 }
	\int_zero:N \l_marine_transpose_row_int
	\prop_clear:N \l_marine_transpose_matrix_prop
	\seq_map_inline:Nn \l_marine_transpose_rows_seq
	{
		\int_incr:N \l_marine_transpose_row_int
		\int_zero:N \l_marine_transpose_col_int
		\seq_set_split:Nnn \l_marine_transpose_arow_seq { & } { ##1 }
		\seq_map_inline:Nn \l_marine_transpose_arow_seq
		{
			\int_incr:N \l_marine_transpose_col_int
			\prop_put:Nxn \l_marine_transpose_matrix_prop
			{
				\int_to_arabic:n { \l_marine_transpose_row_int }
				,
				\int_to_arabic:n { \l_marine_transpose_col_int }
			}
			{ ####1 }
		}
	}
	\tl_clear:N \l_marine_transpose_body_tl
	\int_step_inline:nnnn { 1 } { 1 } { \l_marine_transpose_col_int }
	{
		\int_step_inline:nnnn { 1 } { 1 } { \l_marine_transpose_row_int }
		{
			\tl_put_right:Nx \l_marine_transpose_body_tl
			{
				\prop_item:Nn \l_marine_transpose_matrix_prop { ####1,##1 }
				\int_compare:nF { ####1 = \l_marine_transpose_row_int } { & }
			}
		}
		\tl_put_right:Nn \l_marine_transpose_body_tl { \\ }
	}
	
	% 2. Wrap the output in the environment specified by #1
	\begin{#1}
		\l_marine_transpose_body_tl
	\end{#1}
}

% 3. Generate a variant that expands the second argument (nV)
\cs_generate_variant:Nn \marine_transpose:nn { nV }
\cs_generate_variant:Nn \prop_put:Nnn { Nx }

% 4. Define all the specific transpose environments
\NewEnviron{bmatrixT}{ \marine_transpose:nV { bmatrix } \BODY }
\NewEnviron{pmatrixT}{ \marine_transpose:nV { pmatrix } \BODY }
\NewEnviron{vmatrixT}{ \marine_transpose:nV { vmatrix } \BODY }
\NewEnviron{BmatrixT}{ \marine_transpose:nV { Bmatrix } \BODY } % Braces {}
\NewEnviron{VmatrixT}{ \marine_transpose:nV { Vmatrix } \BODY } % Double bars ||

\ExplSyntaxOff

% ---------------------------------------------------
% GRAPHICS AND TABLES.
% ---------------------------------------------------
% Package for improving the interface with floating objects, such as graphics.
\usepackage{float}
\usepackage{graphicx}
% Packages that offer greater flexibility when using captions for figures, etc.
\usepackage{caption}
\captionsetup{justification=centering}
\usepackage{subcaption}
\usepackage{epstopdf}
% -----------------------
\usepackage{booktabs}
\usepackage{multirow}

% ---------------------------------------------------
% GLOSSARY / ACRONYMS
% ---------------------------------------------------
% Required for the aligned glossary style
\usepackage{longtable}
% The 'acronym' option creates a dedicated list. The 'toc' option includes the glossary in the TOC.
\usepackage[toc,acronym,nomain,nonumberlist]{glossaries}
% No external programs needed to sort the entries. 
\makenoidxglossaries

% Force the acronyms to be bold.
\renewcommand{\glsnamefont}[1]{\textbf{#1}}

% The acronyms definitions will be aligned, and not printed like a standard dictionary (jagged text).
% Create a Custom "Flush Left" Style.
% Tweak the text width parameter 'p{0.XX\textwidth}' accordingly. It tells the second column (the acronyms descriptions) to take up XX% of the page width and wrap nicely like a normal paragraph.
% 1. Define your custom strictly left-justified style
\newglossarystyle{flushleft}{%
	\setglossarystyle{long}% Base it on the existing 'long' style
	\renewenvironment{theglossary}%
	{\begin{longtable}{@{} l p{0.85\textwidth} @{}}}
		{\end{longtable}}%
}
% 2. Tell the package to use the new style
\setglossarystyle{flushleft}

\newacronym{SFU}{SFU}{Solar Flux Unit}

% ---------------------------------------------------
% BIBLIOGRAPHY AND APPENDICES.
% ---------------------------------------------------
% The Backend: biblatex works best with a sorting engine called biber. Explicitly stating backend=biber prevents LaTeX from accidentally trying to use older, broken engines.
\usepackage[
backend=biber,
style=ieee,       	% Pick your citation style (numeric, apa, ieee, authoryear)
maxbibnames=99,   	% Usually, you want all authors in the final list...
minbibnames=1,  	% Ensures no conflict in the bibliography
maxcitenames=2,   	% ...but only 1 or 2 authors when citing inline
mincitenames=1,
uniquelist=false
]{biblatex}

% Use 'family-given' instead of 'last-first' which is technically deprecated in modern LaTeX to avoid hidden warnings.
\DeclareNameAlias{sortname}{family-given}
\DeclareNameAlias{default}{family-given}

% \adddot ensures the period is handled correctly.
\DefineBibliographyStrings{spanish}{%
	andothers = {et al\adddot},
}

\begin{filecontents*}{latex-snippets-bibliography.bib}
	@article{paper_servidia2021,
		author = "P. A. Servidia and M. España",
		title = "On Autonomous Reconfiguration of SAR Satellite Formation Flight With Continuous Control",
		journal = "IEEE Tran. on Aero. and Elec. Systems",
		volume = "57",
		year = "2021",
		month = "April"
	}

\end{filecontents*}
% Link the temporary file created above
\addbibresource{latex-snippets-bibliography.bib} 


% ---------------------------------------------------
% MISC.
% ---------------------------------------------------
% This pack­age sim­pli­fies the in­clu­sion of ex­ter­nal multi-page PDF doc­u­ments in LaTeX doc­u­ments. 
\usepackage{pdfpages}
\setboolean{@twoside}{false}

% ------------------------------------------------------------------------
% SOURCE CODE
% ------------------------------------------------------------------------
\usepackage[dvipsnames,svgnames]{xcolor} 
\usepackage{listings}
\usepackage{caption} % Better handling of code captions
% Algorithms and pseudocode handling.
\usepackage{algorithm}
\usepackage{algpseudocode}
\makeatletter
\renewcommand{\ALG@name}{Algoritmo}
\renewcommand{\listalgorithmname}{Lista de \ALG@name s}
\algrenewcommand\algorithmicensure{\textbf{Salida:}}
\algrenewcommand\algorithmicrequire{\textbf{Entrada:}}
\algnewcommand{\LineComment}[1]{\State \(\triangleright\) #1}
\makeatother

% ------------------------------------------------------------------------
% GLOBAL SETTINGS (Applied to all code)
% ------------------------------------------------------------------------
\lstset{
	inputencoding=utf8,
	extendedchars=true,
	literate={á}{{\'a}}1 {é}{{\'e}}1 {í}{{\'i}}1 {ó}{{\'o}}1 {ú}{{\'u}}1 {ñ}{{\~n}}1 {Á}{{\'A}}1 {É}{{\'E}}1 {Í}{{\'I}}1 {Ó}{{\'O}}1 {Ú}{{\'U}}1 {ü}{{\"u}}1 {Ü}{{\"U}}1,
	basicstyle=\ttfamily\small,
	breaklines=true,			% Enable wrapping
	breakatwhitespace=false,    % Wrap even if there isn't a space
	captionpos=b,
	keepspaces=true,
	showstringspaces=false,
	upquote=true,
	columns=fullflexible,       % Improves character spacing after wrapping
	postbreak=\mbox{\textcolor{red}{$\hookrightarrow$}\space} % Adds a little red arrow at the wrap point.
}

% Spanish Localization
\renewcommand{\lstlistingname}{Código fuente}       	% For the caption
\renewcommand{\lstlistlistingname}{Índice de códigos} 	% For the table of contents

% ------------------------------------------------------------------------
% DEFINE CUSTOM STYLES
% ------------------------------------------------------------------------

% 1. C/C++ Style
\lstdefinestyle{StyleC}{
	language=C++,
	commentstyle=\color{ForestGreen},
	keywordstyle=\color{RoyalBlue},
	stringstyle=\color{BrickRed},
	numberstyle=\tiny\color{gray},
	numbers=left,
	frame=lines,
	backgroundcolor=\color{gray!5}
}

% 2. YAML Style
\lstdefinelanguage{yaml}{
	keywords={true,false,null,y,n},
	keywordstyle=\color{darkgray}\bfseries,
	sensitive=false,
	comment=[l]{\#},
	morecomment=[s]{/*}{*/},
	commentstyle=\color{purple}\ttfamily,
	stringstyle=\color{blue}\ttfamily,
	moredelim=[l][\color{orange}]{\&},
	moredelim=[l][\color{magenta}]{*},
	moredelim=**[il][\color{red}{:}\color{blue}]{:},
	morestring=[b]',
	morestring=[b]",
	frame=lines,
	backgroundcolor=\color{gray!5}
}

% 3. Terminal Style (For console output)
\lstdefinestyle{terminal}{
	backgroundcolor=\color{black},
	basicstyle=\ttfamily\small\color{white},
	frame=lines,
	numbers=none
}